% TEMPLATE-MILC-v3.tex - plantilla maestra UNICA de las guias MILC v3.
%
% La usan las cuatro familias de guias, cada una con su driver:
%   · Grados 6-11          -> scripts/build-guias-g11.py
%   · Bebras               -> scripts/build-guias-bebras.py
%   · Semillero            -> scripts/build-guias-semillero.py
%   · Territorio interior  -> scripts/build-guias-territorio-interior.py
%
% Antes habia tres copias casi identicas (~400 lineas iguales, 6 distintas):
% cada mejora habia que aplicarla tres veces, y en la practica no se hacia
% (el motor de recursos vivia solo en dos de ellas). Lo que varia entre
% familias esta parametrizado en 6 placeholders que cada driver llena
% (se nombran aqui SIN su sintaxis real: el builder reemplaza por texto plano,
% tambien dentro de los comentarios, y un valor multilinea rompe el preambulo):
%
%   HEADER_IZQ        encabezado izquierdo de pagina
%   HEADER_DER        encabezado derecho de pagina
%   PORTADA_ETIQUETA  rotulo sobre el numero grande (GRADO/MOMENTO/MODULO)
%   PORTADA_SERIE     linea de serie de la portada
%   PORTADA_ID        identificador de la guia en la portada
%   PORTADA_CIRCULO   el circulito de grado (°), o vacio si no aplica
%   PDF_TITULO        titulo en texto plano (sin \textbf ni comillas TeX) para
%                     los metadatos del PDF (/Title); lo produce lib_xelatex.texto_pdf
%
% Composicion (auditoria 2026-09-03): las bandas de estacion piden espacio
% (\Needspace*) y prohiben el salto justo despues (\nopagebreak), asi no quedan
% huerfanas al pie; las cajas largas (softbox, drawbox, infoband, puentebox) son
% `breakable`, asi una caja de media pagina ya no arrastra todo a la siguiente
% dejando la anterior al 40 %. Todo texto sobre mostaza o verde va en milcNegro
% (blanco sobre #E5B400 da 1,93:1 y sobre #58A923 2,95:1; ninguno pasa WCAG AA).
% El resto de claves las documenta content/guias/_SCHEMA.md. El script valida
% que no queden placeholders sin reemplazar antes de escribir el archivo final.

% Etiquetado PDF (latex-lab): /Lang, estructura y texto alternativo de las
% figuras, para lectores de pantalla. Debe ir ANTES de \documentclass.
\DocumentMetadata{lang=es-CO,tagging=on}
\documentclass[11pt,letterpaper]{article}
% headheight=24pt: la cabecera derecha lleva dos lineas (identidad de la guia
% y estacion actual). headsep se acorta en la misma medida para que el bloque
% de texto quede exactamente donde estaba con headheight=12pt.
\usepackage[margin=2.0cm,headheight=24pt,headsep=13pt]{geometry}
\usepackage[table]{xcolor}
\usepackage{fontspec}
\usepackage[spanish]{babel}
\usepackage{tikz}
% Librerías TikZ para los diagramas de las guías (bloque `recursos:`):
%  · babel  → babel en español vuelve activos caracteres como «>», y TikZ los usa
%             en sus opciones (>=stealth, ->). Sin esta librería, cualquier
%             diagrama con flechas revienta con «\language@active@arg> has an
%             extra }». Debe ir ANTES de que se dibuje nada.
%  · positioning        → sintaxis «below left=2pt and 3pt».
%  · decorations.pathreplacing → llaves ({brace}) para señalar rangos.
\usetikzlibrary{babel,positioning,decorations.pathreplacing}
\usepackage{graphicx}
% Circuitos: el trabajo de electrónica se hace hoy con micro:bit y sus sensores
% integrados (MakeCode, sin cableado), así que no se necesitan esquemáticos.
% Si algún día entran montajes con protoboard, basta descomentar esta línea:
% los diagramas .tex/.tikz declarados en `recursos:` ya se incrustan solos.
% \usepackage{circuitikz}
\usepackage[most]{tcolorbox}
\usepackage{tabularx}
\usepackage{array}
\usepackage{fancyhdr}
\usepackage{titlesec}
\usepackage[document]{ragged2e}  % bandera a la derecha en todo el cuerpo
\usepackage{enumitem}
\usepackage{microtype}
\usepackage{etoolbox}
\usepackage{needspace}
% hyperref al final del preambulo: metadatos del PDF (titulo, autor, idioma,
% familia), marcadores por estacion (\pdfbookmark) y enlaces sin recuadro.
\usepackage[hidelinks,
  pdftitle={Estética de la liberación — cultura popular y producción ética},
  pdfauthor={PhD. Álvaro Cárdenas Orozco},
  pdfsubject={Tecnología e Informática · Grado 8 · Guía 28 · MILC v3},
  pdflang={es-CO},
  bookmarksopen=true,bookmarksnumbered=false]{hyperref}

% Helvetica Neue no trae la flecha U+2192, asi que cada «→» escrito literalmente
% en el YAML se compilaba a NADA, en silencio (462 flechas en 61 guias: rutas del
% tipo «paleta → tipografia → lockup» salian rotas). Mapeamos el caracter a su
% version matematica, que si tiene glifo. Asi el autor puede seguir escribiendo
% «→» comodamente en el YAML.
\usepackage{newunicodechar}
\newunicodechar{→}{\ensuremath{\rightarrow}}
% Mismo problema con los simbolos de marca y vinetas: el «✓» de «una palomita ✓»
% salia como cuadrito vacio en el PDF de todas las guias que lo usaban.
\usepackage{pifont}
\newunicodechar{✓}{\ding{51}}
\newunicodechar{✗}{\ding{55}}
\newunicodechar{✔}{\ding{52}}
\newunicodechar{•}{\textbullet}
\newunicodechar{≥}{\ensuremath{\geq}}
\newunicodechar{≤}{\ensuremath{\leq}}

\setmainfont{Helvetica Neue}
\setsansfont{Helvetica Neue}
\setlength{\parindent}{0pt}
\setlength{\parskip}{6pt}
% Sin particion de palabras: en 6.º-7.º una palabra cortada por guion al final
% de linea («Comuni-cacion») frena la lectura mucho mas que una linea corta.
\hyphenpenalty=10000
\exhyphenpenalty=10000
\pretolerance=2000
\tolerance=3000
\setlength{\emergencystretch}{3em}
\linespread{1.10}
\renewcommand{\arraystretch}{1.25}
\setlist{leftmargin=5mm,itemsep=1mm,topsep=1mm}
\newcolumntype{Y}{>{\RaggedRight\arraybackslash}X}

\definecolor{milcMagenta}{HTML}{E6007E}
\definecolor{milcVino}{HTML}{5A0038}
\definecolor{milcTurquesa}{HTML}{0093A5}
\definecolor{milcVerde}{HTML}{58A923}
\definecolor{milcMostaza}{HTML}{E5B400}
\definecolor{milcOcre}{HTML}{B86600}
\definecolor{milcNegro}{HTML}{241816}
\definecolor{milcGris}{HTML}{F7F7F4}
\definecolor{milcLinea}{HTML}{E8E2DD}
\definecolor{milcGradePrimary}{HTML}{FF6600}
\definecolor{milcGradeDark}{HTML}{9A3E00}
\definecolor{milcGradeSoft}{HTML}{FFE4D0}
\definecolor{milcGradeAccent}{HTML}{CFE7FF}
\definecolor{milcGradeText}{HTML}{FFFFFF}
\color{milcNegro}

\pagestyle{fancy}
\fancyhf{}
% Cabecera de dos lineas: identidad de la guia arriba y, a la derecha y en
% gris, la estacion en curso (\markright desde \milcestacion). Antes de la
% primera estacion la segunda linea va vacia; el \strut mantiene la altura
% para que la linea de arriba no baile entre paginas.
\lhead{\small Tecnología e Informática · Grado 8\\[-1pt]{\footnotesize\strut}}
\rhead{\small Guía 28 · MILC v3\\[-1pt]{\footnotesize\strut\color{milcNegro!65}\rightmark}}
\cfoot{\small\thepage}
\renewcommand{\headrulewidth}{0pt}

% Sobre mostaza (#E5B400) y verde (#58A923) el texto va en milcNegro; sobre el
% resto de la paleta, en blanco. Se usa dentro de fonttitle/fontupper, donde un
% \color posterior manda sobre coltitle.
\newcommand{\milctintasobre}[1]{%
  \ifboolexpr{test{\ifstrequal{#1}{milcMostaza}} or test{\ifstrequal{#1}{milcVerde}}}%
    {\color{milcNegro}}{\color{white}}}

\newtcolorbox{titlebox}[1]{
  enhanced,colback=#1,colframe=#1,arc=7mm,boxrule=0pt,
  left=7mm,right=7mm,top=3mm,bottom=3mm,
  before skip=8pt,after skip=8pt,
  fontupper=\sffamily\bfseries\Large\color{white}
}

\newtcolorbox{softbox}[3]{
  enhanced,breakable,pad at break=3mm,
  colback=#2,colframe=#1,borderline west={4pt}{0pt}{#1},
  arc=2mm,boxrule=.4pt,left=4mm,right=4mm,top=3mm,bottom=3mm,
  before skip=6pt,after skip=8pt,title={#3},coltitle=white,
  colbacktitle=#1,fonttitle=\sffamily\bfseries\milctintasobre{#1},fontupper=\RaggedRight,
  attach boxed title to top left={xshift=3mm,yshift=-2mm},
  boxed title style={arc=2mm, boxrule=0pt}
}

\newtcolorbox{drawbox}[2]{
  enhanced,breakable,pad at break=4mm,
  colback=white,colframe=#1,arc=4mm,boxrule=1pt,
  left=5mm,right=5mm,top=4mm,bottom=4mm,before skip=8pt,after skip=8pt,
  title={#2},coltitle=white,colbacktitle=#1,fonttitle=\sffamily\bfseries\milctintasobre{#1},
  fontupper=\RaggedRight,
  attach boxed title to top left={xshift=4mm,yshift=-2mm},
  boxed title style={arc=3mm, boxrule=0pt}
}

\newtcolorbox{infoband}[1]{
  enhanced,breakable,pad at break=2mm,colback=#1!8,colframe=#1!50,arc=2mm,boxrule=.5pt,
  left=4mm,right=4mm,top=2mm,bottom=2mm,before skip=4pt,after skip=4pt,
  fontupper=\small\RaggedRight
}

% Puente narrativo entre secciones (contrato editorial Conéctate).
% Da continuidad: "Acabas de X. Ahora vas a Y porque Z."
% Visualmente: banda izquierda en color del grado, texto en itálica pequeña.
\newtcolorbox{puentebox}{
  enhanced,breakable,pad at break=2mm,colback=milcGradeSoft,colframe=milcGradeDark,
  borderline west={3pt}{0pt}{milcGradeDark},
  arc=0mm,boxrule=0pt,left=5mm,right=4mm,top=2.5mm,bottom=2.5mm,
  before skip=4pt,after skip=8pt,
  fontupper=\itshape\small\color{milcNegro}\RaggedRight
}
% La flecha va en modo matematico a proposito: Helvetica Neue NO tiene el glifo
% U+2192, asi que \textrightarrow se compilaba a NADA («Missing character: There
% is no → in font Helvetica Neue») y el puente salia sin flecha. En modo
% matematico la flecha viene de la fuente matematica y si se dibuja.
\newcommand{\puente}[1]{%
  \begin{puentebox}%
    \textcolor{milcGradeDark}{$\rightarrow$}\hspace{2mm}#1%
  \end{puentebox}%
}

\newcommand{\linead}[1]{\noindent\color{milcLinea}\rule{#1}{0.5pt}\color{milcNegro}}
\newcommand{\respuesta}[1]{\par\vspace{2mm}\foreach \n in {1,...,#1}{\noindent\color{milcLinea}\rule{\linewidth}{0.5pt}\par\vspace{5mm}}\color{milcNegro}}
\newcommand{\checkbox}{\textcolor{milcGradeDark}{\fbox{\phantom{x}}}\hspace{5pt}}

% ─── Listas aireadas para las guías socioemocionales ────────────────────
% Componer las verificaciones como «\checkbox texto\\» deja el texto que salta
% de línea debajo del cuadrito y apelmaza el bloque. Estos entornos alinean la
% continuación y airean los ítems. Los usa Territorio Interior; cualquier
% familia puede hacerlo.
\newlist{checklista}{itemize}{1}
\setlist[checklista]{
  label=\textcolor{milcGradeDark}{\fbox{\phantom{x}}},
  leftmargin=9mm, labelsep=3.2mm, itemsep=2.4mm,
  topsep=2.5mm, parsep=0pt, partopsep=0pt, before=\RaggedRight
}
\newlist{pasoslista}{enumerate}{1}
\setlist[pasoslista]{
  label=\textcolor{milcGradeDark}{\sffamily\bfseries\arabic*.},
  leftmargin=8mm, labelsep=3mm, itemsep=2.8mm,
  topsep=1mm, parsep=0pt, partopsep=0pt, before=\RaggedRight
}

\newcommand{\phasecard}[4]{
\begin{tcolorbox}[enhanced,colback=#3,colframe=#2,arc=3mm,boxrule=.5pt,height=3.05cm,valign=center,halign=center,left=1.5mm,right=1.5mm,top=2mm,bottom=2mm]
{\sffamily\bfseries\fontsize{22}{24}\selectfont\textcolor{#2}{#1}}\par
{\sffamily\bfseries\small #4}
\end{tcolorbox}
}

% ── Piezas del rediseno 2026 (legibilidad 6.º-11.º) ───────────────────
% Banda de estacion: reemplaza el titlebox macizo. Titulo a la izquierda,
% dato practico (tiempo/modalidad) a la derecha, en una sola linea.
% Nunca huerfana: pide 9 lineas libres (o salta de pagina rellenando con
% \vfill) y prohibe el salto justo despues de la banda. Nueve y no seis:
% la caja partible que sigue necesita ~80pt (titulo adosado, relleno y dos
% lineas) para arrancar en la misma pagina; con menos, tcolorbox la manda
% entera a la siguiente y la banda se queda sola. El penalty 10000 va
% en `after`, ANTES del espacio posterior: un `after skip` (aunque sea 0pt)
% es un \vskip tras la caja, y ese glue es un punto de corte legal que un
% \nopagebreak escrito despues de \end{tcolorbox} ya no cierra. Cada banda
% deja marcador PDF y actualiza la cabecera derecha.
\newcounter{milcestacion}
\newcommand{\milcestacion}[2]{%
\Needspace*{9\baselineskip}%
\stepcounter{milcestacion}%
\pdfbookmark[1]{#2}{milcestacion\themilcestacion}%
\markright{#2}%
\begin{tcolorbox}[enhanced,colback=#1,colframe=#1,arc=2mm,boxrule=0pt,
  left=5mm,right=5mm,top=2.6mm,bottom=2.6mm,before skip=11pt,
  after={\par\nopagebreak\vskip7pt}]
{\sffamily\bfseries\fontsize{15}{18}\selectfont\milctintasobre{#1}#2}
\end{tcolorbox}}

% Tarjeta de paso numerada: sustituye las tablas «Pilar / Aplicacion», que
% eran formato de consulta docente, no de estudiante. El numero va en un
% circulo sobre el borde para que la secuencia se lea de un vistazo.
\newcommand{\milcpaso}[3]{%
\Needspace{4\baselineskip}%
\begin{tcolorbox}[enhanced,colback=white,colframe=milcLinea,arc=1.5mm,boxrule=.9pt,
  left=14mm,right=4mm,top=3mm,bottom=3mm,before skip=4pt,after skip=4pt,
  fontupper=\RaggedRight,
  overlay={\node[anchor=north west,circle,fill=#2,text=white,inner sep=0pt,
    minimum size=7.4mm,font=\sffamily\bfseries\small]
    at ([xshift=3.6mm,yshift=-3.2mm]frame.north west) {#1};}]
#3
\end{tcolorbox}}

% Casilla marcable de tamano util con lapiz (la anterior era \fbox{\phantom{x}},
% de alto variable segun el texto que la seguia).
\newcommand{\milcmarca}{\framebox[3.8mm]{\rule{0pt}{3.4mm}}}

% Fila marcable de lista suelta (checks de la estacion 1).
\newcommand{\milccheckrow}[1]{%
\par\vspace{1.8mm}\noindent
\makebox[6.4mm][l]{\raisebox{-.55mm}{\textcolor{milcNegro}{\milcmarca}}}%
\parbox[t]{\dimexpr\linewidth-6.4mm\relax}{\RaggedRight #1}\par}

% Celda marcable dentro de la rubrica (tabla de autoevaluacion). Misma
% sangria francesa, pero pensada para vivir dentro de una columna X.
\newcommand{\milcopcion}[1]{%
\makebox[5.2mm][l]{\raisebox{-.55mm}{\textcolor{milcNegro}{\milcmarca}}}%
\parbox[t]{\dimexpr\linewidth-5.2mm\relax}{\RaggedRight #1}}

% ── Recursos visuales de la guía ──────────────────────────────────────
% Declarados en el bloque `recursos:` del YAML y emitidos por el builder.
% \guiaFigura   → imagen rasterizada o PDF (foto, carta celeste, montaje).
% \guiaDiagrama → fuente TikZ/circuitikz (.tex) incrustada como vector:
%                 es la vía para circuitos y esquemáticos editables.
% [#1] = texto alternativo (alt) · #2 = ruta relativa al .tex · #3 = pie
% El alt llega al PDF etiquetado (/Alt) y lo lee el lector de pantalla.
\newcommand{\guiaFigura}[3][]{%
\begin{center}
\includegraphics[width=0.88\linewidth,height=0.38\textheight,keepaspectratio,alt={#1}]{#2}\\[1.5mm]
{\footnotesize\itshape\textcolor{milcNegro!75}{#3}}
\end{center}
}

\newcommand{\guiaDiagrama}[2]{%
\begin{center}
\input{#1}\\[1.5mm]
{\footnotesize\itshape\textcolor{milcNegro!75}{#2}}
\end{center}
}

\begin{document}
\thispagestyle{empty}

% ── Portada ───────────────────────────────────────────────────────────
% Antes: un rectangulo de color a pagina completa. Se veia bien en pantalla,
% pero en la fotocopiadora del colegio se comia el toner y salia gris sucio.
% Ahora la banda ocupa el tercio superior y el resto va en blanco.
\begin{tikzpicture}[remember picture,overlay]
  \path[fill=milcGradePrimary]
    (current page.north west) rectangle ([yshift=-10.6cm]current page.north east);

  \node[anchor=north west,text=milcGradeText] at ([xshift=2.0cm,yshift=-1.55cm]current page.north west)
    {\sffamily\bfseries\fontsize{12}{14}\selectfont GRADO};
  \node[anchor=north west,text=milcGradeText,opacity=.85] at ([xshift=2.0cm,yshift=-2.35cm]current page.north west)
    {\sffamily\bfseries\fontsize{8.5}{10}\selectfont SERIE GUÍAS · TECNOLOGÍA E INFORMÁTICA};
  \node[anchor=north east,text=milcGradeText,opacity=.85,align=right] at ([xshift=-2.0cm,yshift=-1.55cm]current page.north east)
    {\sffamily\bfseries\fontsize{9}{12}\selectfont I.E. Sor María Juliana\\Cartago, Valle del Cauca};

  \node[anchor=west,text=milcGradeText] (gradonum) at ([xshift=1.85cm,yshift=-7.1cm]current page.north west)
    {\sffamily\bfseries\fontsize{112}{112}\selectfont 8};
  % El circulito de grado (°) solo aplica a las guias de grado («11°»); en
  % Bebras y Semillero el numero grande es un momento/modulo, no un grado.
  % Se posiciona relativo al nodo `gradonum`, no en coordenadas absolutas.
  \draw[milcGradeText,line width=1.6pt]
    ([xshift=2.0mm,yshift=-4.5mm]gradonum.north east) circle (.15cm);

  \node[anchor=west,text=milcGradeText,text width=11.4cm,align=left] at ([xshift=1.5cm,yshift=0.5cm]gradonum.east)
    {
     {\sffamily\bfseries\fontsize{9}{11}\selectfont\MakeUppercase{Período 3 · Multimedia, narrativa y ciberseguridad}}\\[3mm]
     {\sffamily\bfseries\fontsize{21}{25}\selectfont\setlength{\baselineskip}{27pt}Estética de la liberación ---\\[4pt]cultura popular\\[4pt]y producción ética}\\[3.5mm]
     {\sffamily\bfseries\fontsize{15}{18}\selectfont Guía 28-8-TIC}
    };
\end{tikzpicture}

\vspace*{9.4cm}
\vfill

{\sffamily\bfseries\fontsize{8}{10}\selectfont\textcolor{milcGradeDark}{LO QUE TE LLEVAS HOY}}

\begin{tcolorbox}[enhanced,colback=white,colframe=milcNegro,arc=2mm,boxrule=1.6pt,
  left=5mm,right=5mm,top=4mm,bottom=4mm,before skip=3pt,after skip=12pt,
  fontupper=\RaggedRight\fontsize{11}{15.5}\selectfont]
Una pieza multimedia sobre un elemento de la cultura popular de tu barrio, del Valle o del Pacífico: afiche, video de un minuto, audio corto o infografía. La acompañan el análisis de una pieza comercial sobre el mismo tema, marcada con los cinco errores que reproducen jerarquías, y cuatro párrafos honestos sobre la tuya: a quién amplifica, a quién dignifica, qué dejaste fuera y qué cambiarías si pudieras consultar a la persona retratada.
\end{tcolorbox}

\begin{tcolorbox}[enhanced,colback=milcGris,colframe=milcGris,arc=2mm,boxrule=0pt,
  left=5mm,right=5mm,top=3.5mm,bottom=3.5mm,after skip=12pt]
{\sffamily\bfseries\fontsize{8}{10}\selectfont\textcolor{milcNegro!55}{ESTA GUÍA ES DE}}\\[3mm]
\begin{tabularx}{\linewidth}{X p{3.2cm} p{3.2cm}}
\linead{\linewidth} & \linead{3.0cm} & \linead{3.0cm}\\[-1mm]
{\fontsize{8.5}{10}\selectfont\textcolor{milcNegro!55}{Nombre completo}} &
{\fontsize{8.5}{10}\selectfont\textcolor{milcNegro!55}{Grupo}} &
{\fontsize{8.5}{10}\selectfont\textcolor{milcNegro!55}{Fecha}}\\
\end{tabularx}
\end{tcolorbox}

\vfill

\noindent\textcolor{milcLinea}{\rule{\linewidth}{1pt}}\\[2mm]
{\sffamily\bfseries\fontsize{9.5}{12}\selectfont Autor: Dr. Álvaro Cárdenas Orozco}\\[1mm]
{\sffamily\fontsize{8.5}{11}\selectfont\textcolor{milcNegro!55}{Metodología MILC · apertura ancestral · pensamiento computacional · triángulo Dussel–estoicismo–Floridi}}

\newpage

\section*{Apertura: Estética de la liberación --- cultura popular y producción ética}

\begin{softbox}{milcGradeDark}{milcGradeSoft}{Saber ancestral}
Alguien nos enseñó a mirar los cuerpos así. La antropóloga Zandra Pedraza revisó lo que se leyó en Colombia entre 1830 y 1990: cartillas escolares, manuales de higiene, textos de educación popular y la revista Cromos. Encontró que el país fue aprendiendo, página por página, qué cuerpo se consideraba sano, decente y digno de respeto, y cuál no. Ese aprendizaje pasó por la escuela. La clase de higiene y la de urbanidad enseñaban a pararse, a comer, a vestirse y, sin decirlo, a avergonzarse. El cuerpo del campesino, del negro y del indígena aparecía muchas veces del lado del atraso. Por eso la burla no es un chiste espontáneo: es una lección vieja que alguien repite sin saber de dónde la sacó. La \textbf{cara de exclusión} señala a la escuela misma: fue el aparato que enseñó cuál cuerpo era decente. Esta guía no puede presentarse como ajena a esa historia. Hoy vas a producir una pieza, y toda pieza enseña a mirar.\par\smallskip{\footnotesize\textcolor{milcNegro!70}{\textbf{Fuente.} Pedraza Gómez, Z. (1999). \emph{En cuerpo y alma: visiones del progreso y de la felicidad}. Universidad de los Andes, Departamento de Antropología.}}
\end{softbox}

\begin{softbox}{milcTurquesa}{milcGris}{Saber contemporáneo}
La \textbf{estética de la liberación} es la propuesta de Enrique Dussel para juzgar una pieza por tres preguntas. ¿A quién amplifica? ¿A quién dignifica? ¿A quién deja fuera? De ahí salen cuatro principios de trabajo. \textbf{Centrar la voz del sur}: el sujeto es alguien del barrio, del Valle o del Pacífico, no un personaje genérico. \textbf{Dignificar a quien aparece}: con nombre, oficio y contexto, no como tipo folclórico. \textbf{Mostrar las contradicciones}: la cultura popular tiene riqueza y también dificultad, y ocultar la segunda es publicidad. \textbf{Producir para la comunidad}: la pieza debe poder mostrarse a la persona retratada, y esa persona debe reconocerse sin vergüenza. Ninguna pieza es neutral. Toda pieza reparte dignidad de algún modo, y la única diferencia es si quien la hizo sabía lo que estaba repartiendo.
\end{softbox}

\begin{softbox}{milcMostaza}{milcGris}{Pregunta-puente}
\textit{Si Colombia aprendió a mirar cuerpos leyendo cartillas y revistas, tu pieza también va a enseñarle a mirar a alguien. La pregunta es sencilla y difícil: ¿qué le va a enseñar?}
\end{softbox}

\begin{softbox}{milcVerde}{milcGris}{Saber y saber hacer}
Al terminar podrás: (1) \textbf{identificar} elementos de la cultura popular de tu entorno que la publicidad casi nunca representa con respeto; (2) \textbf{analizar} una pieza comercial con los cuatro principios y los cinco errores que reproducen jerarquías; (3) \textbf{crear} tu propia pieza y sostener por escrito a quién amplifica, a quién dignifica y qué dejaste fuera.
\end{softbox}



\small
\begin{tabularx}{\textwidth}{>{\bfseries}p{3.0cm}Y}
\rowcolor{milcGradePrimary!12}Autor & Dr. Álvaro Cárdenas Orozco\\
DBA & Producción multimedia ética y comportamiento responsable en entornos digitales (MEN, T\&I)\\
Referentes & Dussel (estética de la liberación) · Ley 1336 · Floridi (infoesfera)\\
Producto final & Una pieza multimedia sobre un elemento de la cultura popular de tu barrio, del Valle o del Pacífico: afiche, video de un minuto, audio corto o infografía. La acompañan el análisis de una pieza comercial sobre el mismo tema, marcada con los cinco errores que reproducen jerarquías, y cuatro párrafos honestos sobre la tuya: a quién amplifica, a quién dignifica, qué dejaste fuera y qué cambiarías si pudieras consultar a la persona retratada.\\
\end{tabularx}
\normalsize

\puente{En las sesiones 1 a 5 aprendiste a producir imagen, relato y audio. En la 6 y la 7 trabajaste la responsabilidad frente al daño. Hoy le agregas la pregunta política: toda pieza toma partido. La sesión 9 abre el proyecto integrador y la 10 lo sustenta.}

\milcestacion{milcGradeDark}{Ruta MILC: así vamos a aprender}
\begin{tabularx}{\textwidth}{>{\centering\arraybackslash}X>{\centering\arraybackslash}X>{\centering\arraybackslash}X>{\centering\arraybackslash}X}
\phasecard{1}{milcVerde}{milcGris}{Escucha\\[-1mm]\small Observo, escucho y pregunto.} &
\phasecard{2}{milcTurquesa}{milcGris}{Entender\\[-1mm]\small Sistematizo y ordeno.} &
\phasecard{3}{milcMagenta}{milcGris}{Hacer\\[-1mm]\small Construyo y aplico.} &
\phasecard{4}{milcVino}{milcGris}{Cierre\\[-1mm]\small Evalúo y reflexiono.}
\end{tabularx}


\puente{Antes de la teoría vas a hacer un inventario corto. Cinco elementos de la cultura popular de tu barrio o de tu región que creas que la publicidad no representa bien. De ahí sale el tema de tu pieza.}

\milcestacion{milcVerde}{Estación 1 de 3 · Escucha}

\begin{softbox}{milcVerde}{milcGris}{Escena para escuchar}
\textbf{Actividad 1 · IDENTIFICA --- Inventario de cultura popular} (15 min · individual). (1) Anota cinco elementos de la cultura popular de tu barrio, del Valle o del Pacífico. Pueden ser oficios, comidas, músicas, fiestas o personas concretas del sector. (2) Al lado de cada uno, escribe en pocas palabras cómo lo suele mostrar la publicidad. (3) Marca los que crees que salen mal parados y explica en una línea por qué. (4) Elige uno para tu pieza y escribe a quién concreto podrías consultar sobre él.
\end{softbox}

{\sffamily\bfseries\fontsize{8}{10}\selectfont\textcolor{milcNegro!55}{MARCA CUANDO LO HAYAS HECHO}}
\milccheckrow{Anoté cinco elementos con cómo los muestra la publicidad.}
\milccheckrow{Marqué los que salen mal parados y dije por qué.}
\milccheckrow{Elegí uno y nombré a alguien concreto a quien consultar.}

\begin{infoband}{milcVerde}


\textbf{Lo que va al cuaderno (Actividad 1 · IDENTIFICA).} \textbf{Título:} «Actividad 1 --- Inventario de cultura popular». \textbf{Formato:} lista de cinco con dos columnas (elemento / cómo lo muestra la publicidad), más el elegido y el nombre de la persona a consultar. \textbf{Extensión:} un tercio de página. \textbf{Sabes que terminaste cuando} tu elección tiene una razón escrita y una persona detrás.
\end{infoband}

\puente{Ya tienes el tema. Ahora vas a tomar una pieza comercial sobre algo parecido y a pasarla por los cuatro principios y los cinco errores, para ver el mecanismo por dentro.}

\milcestacion{milcTurquesa}{Estación 2 de 3 · Entender}

\begin{softbox}{milcTurquesa}{milcGris}{Lo que vas a entender}
\textbf{Actividad 2 · ANALIZA --- Una pieza comercial por dentro} (30 min · parejas). (1) Con tu pareja, elijan una publicidad o una pieza de medios masivos sobre cultura popular colombiana. (2) Pásenla por los cuatro principios y anoten cuáles cumple y cuáles no. (3) Marquen cuál de los cinco errores aparece y copien la parte concreta de la pieza donde se ve. (4) Escriban en dos líneas qué le enseña esa pieza a quien la mira sobre la gente que retrata.
\end{softbox}

\milcpaso{1}{milcTurquesa}{\textbf{Los cuatro principios.} \textbf{Centrar la voz del sur}: quien aparece es de aquí y se le nota de dónde. \textbf{Dignificar}: aparece con nombre, oficio y contexto, no como tipo. \textbf{Mostrar las contradicciones}: la riqueza y la dificultad, las dos. \textbf{Producir para la comunidad}: la pieza se le puede mostrar a la persona retratada.}
\milcpaso{2}{milcTurquesa}{\textbf{Los cinco errores.} \textbf{Caricaturización}: el rasgo se exagera hasta el chiste. \textbf{Romantización}: solo lo bonito, sin las dificultades reales. \textbf{Exotismo}: la persona aparece como rareza para quien mira desde fuera. \textbf{Paternalismo}: aparece como víctima que espera salvación, no como alguien que decide. \textbf{Extractivismo cultural}: se toman música, símbolos o palabras sin nombrar de dónde vienen ni devolver nada.}
\milcpaso{3}{milcTurquesa}{\textbf{Cómo se reconoce el exotismo.} Pregúntate para quién está hecha la pieza. Si el placer de mirarla depende de que quien mira no sea del lugar, ahí está. La cultura popular no necesita volverse extraña para valer.}
\milcpaso{4}{milcTurquesa}{\textbf{La pregunta que cierra el análisis.} ¿La persona retratada se reconocería en esta pieza con orgullo? Si la respuesta es sí, la pieza dignifica. Si es no, o si dudas, la pieza está repitiendo la lección vieja de las cartillas.}

\begin{drawbox}{milcTurquesa}{Estética de la liberación, en una mirada}
\textbf{3 preguntas:} a quién amplifica / a quién dignifica / a quién deja fuera. \\[1mm] \textbf{4 principios:} voz del sur / dignidad / contradicciones / para la comunidad. \\[1mm] \textbf{5 errores:} caricatura / romantización / exotismo / paternalismo / extractivismo. \\[1mm] \textbf{Prueba:} ¿se reconocería con orgullo?
\end{drawbox}

\begin{softbox}{milcMostaza}{milcGris}{Errores comunes y su corrección}
(1) \textbf{Creer que una pieza puede ser neutral}: toda pieza reparte dignidad, se sepa o no. (2) \textbf{Confundir dignificar con embellecer}: quitar las dificultades es romantizar, no respetar. (3) \textbf{Usar a una persona sin nombrarla}: sin nombre y sin oficio, queda convertida en tipo. (4) \textbf{Tomar música o símbolos sin acreditar}: eso es extractivismo, aunque la intención sea buena. (5) \textbf{No mostrarle la pieza a nadie del lugar}: es el único chequeo que de verdad sirve.
\end{softbox}

\begin{infoband}{milcTurquesa}


\textbf{Lo que va al cuaderno (Actividad 2 · ANALIZA).} \textbf{Título:} «Actividad 2 --- Una pieza comercial por dentro». \textbf{Formato:} los cuatro principios con cumple o no cumple, el error identificado con la parte concreta donde se ve, y dos líneas sobre qué enseña la pieza. \textbf{Extensión:} media página. \textbf{Sabes que terminaste cuando} puedes señalar el error en un fragmento concreto, no en general.
\end{infoband}

\puente{Tienes el criterio. Toca producir la tuya y sostenerla por escrito, con la parte incómoda incluida: qué dejaste fuera y por qué.}

\milcestacion{milcMagenta}{Estación 3 de 3 · Hacer}

\begin{softbox}{milcMagenta}{milcGris}{Cómo vas a construir tu producto}
\textbf{Actividad 3 · CREA --- Tu pieza y sus cuatro párrafos} (25 min · individual). (1) Produce tu pieza sobre el elemento que elegiste, con el formato que le convenga. (2) Asegúrate de que quien aparece tenga nombre, oficio y contexto. (3) Incluye una contradicción, algo que no sea solo lo bonito del tema. (4) Escribe los cuatro párrafos: a quién amplifica, a quién dignifica, qué dejaste fuera y qué cambiarías si pudieras consultarla. (5) Revisa tu pieza contra los cinco errores y anota cuál estuviste a punto de cometer.
\end{softbox}

\milcpaso{1}{milcMagenta}{\textbf{Tu entrega tiene tres piezas.} La pieza multimedia. El análisis de la pieza comercial de la Actividad 2. Los cuatro párrafos sobre la tuya.}
\milcpaso{2}{milcMagenta}{\textbf{El formato sigue al tema.} Un oficio pide afiche o retrato. Una fiesta pide video corto. Una música pide audio. Una comida pide infografía. Elegir bien el formato ya es parte del respeto.}
\milcpaso{3}{milcMagenta}{\textbf{El párrafo difícil es el tercero.} Qué dejaste fuera. Toda pieza recorta, y decir qué quedó por fuera es más honesto que fingir que cabía todo. Ahí se ve si entendiste la idea.}
\milcpaso{4}{milcMagenta}{\textbf{La prueba final no es del profesor.} Es imaginar a la persona retratada viendo tu pieza. Si te incomoda esa imagen, la pieza todavía no está.}

\begin{drawbox}{milcMagenta}{Antes de entregar}
\checkbox Pieza sobre el elemento elegido en el inventario. \\[1mm] \checkbox Quien aparece tiene nombre, oficio y contexto. \\[1mm] \checkbox La pieza muestra una contradicción, no solo lo bonito. \\[1mm] \checkbox Análisis de la pieza comercial con el error señalado en un fragmento. \\[1mm] \checkbox Los cuatro párrafos, incluido qué dejaste fuera. \\[1mm] \checkbox Revisión contra los cinco errores.
\end{drawbox}

\begin{softbox}{milcMagenta}{milcGris}{Plantilla de trabajo}
\textbf{Tema.} \textit{Elemento del inventario y por qué.} \\[1mm] \textbf{Quien aparece.} \textit{Nombre, oficio, contexto.} \\[1mm] \textbf{Pieza.} \textit{Formato y herramienta.} \\[1mm] \textbf{Cuatro párrafos.} \textit{Amplifica, dignifica, dejé fuera, cambiaría.} \\[1mm] \textbf{Comparación.} \textit{La pieza comercial y su error.}
\end{softbox}

\begin{infoband}{milcMagenta}


\textbf{Lo que va al cuaderno (Actividad 3 · CREA).} \textbf{Título:} «Actividad 3 --- Tu pieza y sus cuatro párrafos». \textbf{Formato:} la referencia a la pieza con su formato, los cuatro párrafos y la nota sobre el error que estuviste a punto de cometer. \textbf{Extensión:} media página. \textbf{Sabes que terminaste cuando} podrías mostrarle la pieza a la persona retratada sin incomodidad.
\end{infoband}

\puente{Lo que entregas hoy: la pieza, el análisis de la pieza comercial y tus cuatro párrafos. La prueba de calidad: la persona retratada podría verla y reconocerse sin vergüenza.}

\milcestacion{milcMostaza}{Producto de la sesión}
\textbf{Pieza multimedia + análisis de una pieza comercial + cuatro párrafos honestos}.
\begin{softbox}{milcMostaza}{milcGris}{Criterios de calidad}
(1) La pieza trata un elemento del inventario, de tu barrio o de tu región. (2) Quien aparece tiene nombre, oficio y contexto, no es un tipo genérico. (3) La pieza muestra una contradicción y no solo lo bonito del tema. (4) El análisis de la pieza comercial señala el error en un fragmento concreto. (5) Los cuatro párrafos dicen qué quedó fuera y qué cambiarías tras consultar.
\end{softbox}

\puente{Antes de cerrar, mira lo que hiciste desde cinco lados. Producir sin preguntarse a quién dignifica una pieza no la vuelve neutral: la vuelve repetición de lo que ya estaba.}

\milcestacion{milcVino}{Evaluación liberadora · cinco dimensiones}

\begin{softbox}{milcVino}{milcGris}{Sentido de la evaluación}
Evaluar no es solo revisar una nota. Es reconocer qué aprendí, qué cambió en mí, qué emociones manejé, cómo me relaciono con la ciudad y la comunidad, y qué le dejaré a quien viene detrás.
\end{softbox}

\textbf{Mi reflexión liberadora en cinco dimensiones.} Completa cada frase iniciada:

\small
\begin{tabularx}{\textwidth}{>{\bfseries\RaggedRight\arraybackslash}p{3.4cm}Y>{\RaggedRight\arraybackslash}p{4.6cm}}
\rowcolor{milcVino}\color{white}Dimensión & \color{white}\bfseries Mi respuesta (frame starter) & \color{white}\bfseries Algo que descubrí\\
1. Personal & Lo que más cambió en mí fue \linead{6.5cm} & \linead{4.2cm}\\
2. Emocional & La emoción más fuerte fue \linead{2.5cm} y la regulé \linead{3cm} & \linead{4.2cm}\\
3. Ciudadana & En democracia esto importa porque \linead{6cm} & \linead{4.2cm}\\
4. Local (Cartago) & En mi barrio sería útil para \linead{6.5cm} & \linead{4.2cm}\\
5. Intergeneracional & A mi abuela/abuelo le diría \linead{6.5cm} & \linead{4.2cm}\\
\end{tabularx}
\normalsize

\begin{infoband}{milcVino}
\textbf{Para la próxima sustentación.} Vuelve a esta tabla en seis meses. Verás que las respuestas cambian — no porque lo aprendido cambie, sino porque tú cambias. La evaluación liberadora es una conversación contigo a través del tiempo.
\end{infoband}


\puente{Para terminar, tres ideas sobre representar a alguien, contadas con palabras de esta guía: Dussel sobre lo que siempre queda por fuera de una mirada, Marco Aurelio sobre hacer en vez de declarar, Floridi sobre el miedo a lo que no entendemos.}

\milcestacion{milcGradeDark}{Tres ideas para cerrar · Dussel, un estoico y Floridi}

\begin{softbox}{milcVino}{milcGris}{Enrique Dussel · lente del nosotros}
\textit{Toda persona y todo pueblo está siempre más allá del sistema desde el que otro los mira.}

{\footnotesize\itshape\color{milcNegro!70}--- idea tomada de Enrique Dussel, contada con palabras de esta guía. No es una cita textual.}

Tu pieza recorta, y quien aparece en ella es más de lo que cabe en el encuadre. Por eso el párrafo de «qué dejé fuera» no es un trámite: es reconocer que la persona no termina donde termina tu foto.

\textbf{¿Qué parte de esta persona no cupo en mi pieza, y por qué elegí dejarla fuera?}
\end{softbox}
\linead{\linewidth}\\[1mm]
\linead{\linewidth}

\begin{softbox}{milcMostaza}{milcGris}{Estoicismo · Marco Aurelio · cuidado interior}
\textit{Deja de discutir cómo debería ser una persona buena y procura serlo de verdad.}

{\footnotesize\itshape\color{milcNegro!70}--- idea tomada de Marco Aurelio, contada con palabras de esta guía. No es una cita textual.}

Una pieza no se vuelve respetuosa porque el análisis diga que lo es. Se vuelve respetuosa en decisiones concretas: el nombre que escribes, el plano que eliges, la frase que no pones. Lo demás es declaración.

\textbf{¿En qué decisión concreta de mi pieza se nota el respeto, sin que yo tenga que decirlo?}
\end{softbox}
\linead{\linewidth}\\[1mm]
\linead{\linewidth}

\begin{softbox}{milcTurquesa}{milcGris}{Luciano Floridi y la Onlife Initiative · lente de la infoesfera}
\textit{Tememos y rechazamos aquello a lo que no logramos darle sentido.}

{\footnotesize\itshape\color{milcNegro!70}--- idea tomada de Luciano Floridi y la Onlife Initiative, contada con palabras de esta guía. No es una cita textual.}

La caricatura y el exotismo son atajos para lo que no se entendió. Cuando alguien no conoce un oficio, lo dibuja gracioso o lo dibuja exótico. Conocer antes de producir es lo que evita las dos salidas.

\textbf{¿Qué de lo que retraté conozco de verdad, y qué estoy suponiendo desde afuera?}
\end{softbox}
\linead{\linewidth}\\[1mm]
\linead{\linewidth}

\begin{infoband}{milcNegro}
\textbf{Nota para el docente.} Las tres ideas están contadas con palabras de esta guía a partir del pensamiento de cada autor; no son citas textuales. De 9.º en adelante se trabajan con la obra y su referencia.
\end{infoband}

\milcestacion{milcVino}{Mi autoevaluación · marca una casilla por fila}

{\small
\setlength{\extrarowheight}{2.2pt}
\begin{tabularx}{\textwidth}{>{\RaggedRight\arraybackslash\bfseries}p{3.0cm}YYY}
\rowcolor{milcVino}\color{white}\bfseries Criterio & \color{white}\bfseries Lo logré & \color{white}\bfseries En proceso & \color{white}\bfseries Necesito apoyo\\[.6mm]
Comprensión del tema & \milcopcion{Explico las ideas con mis palabras.} & \milcopcion{Reconozco las ideas, falta precisión.} & \milcopcion{Necesito repasar lo básico.}\\[.8mm]
\rowcolor{milcGris}Pensamiento computacional & \milcopcion{Usé los pilares al armar mi producto.} & \milcopcion{Usé algunos pilares.} & \milcopcion{Necesito releer la estación 2.}\\[.8mm]
Aplicación y producto & \milcopcion{Mi producto cumple los criterios.} & \milcopcion{Está armado, con ajustes pendientes.} & \milcopcion{Aún está en construcción.}\\[.8mm]
\rowcolor{milcGris}Reflexión 5 dimensiones & \milcopcion{Respondí cada una con honestidad.} & \milcopcion{Respondí algunas con detalle.} & \milcopcion{Mis respuestas son muy generales.}\\[.8mm]
Triángulo de pensamiento & \milcopcion{Conecté las 3 voces con mi trabajo.} & \milcopcion{Conecté al menos una.} & \milcopcion{No las relaciono con mi proceso.}\\[.8mm]
\end{tabularx}}

\vspace{3mm}
\textbf{Hoy entendí que:}\\[1mm]
\linead{\linewidth}\\[1mm]
\linead{\linewidth}

\textbf{Mi compromiso para la próxima sesión es:}\\[1mm]
\linead{\linewidth}\\[1mm]
\linead{\linewidth}

\begin{softbox}{milcMostaza}{milcGris}{Cita de despedida · Marco Aurelio}
\textit{``El obstáculo en el camino se vuelve el camino. Lo que estorba al camino se vuelve camino.''}\\[1mm]
Cada error de hoy, bien observado, se convierte en pieza del camino que pisarás mañana.
\end{softbox}

\small
\begin{tabularx}{\textwidth}{>{\bfseries}p{3.6cm}Y}
\rowcolor{milcGradeDark!12}Nombre completo: & \linead{10cm}\\
Fecha: & \linead{4cm} \quad Curso/grupo: \linead{4cm}\\
Mi firma: & \linead{9cm}\\
\end{tabularx}
\normalsize

\begin{infoband}{milcGradeDark}
\textbf{Para la próxima sesión.} Llega con la pieza, el análisis y los cuatro párrafos. Si puedes, muéstrale la pieza a la persona que retrataste y anota qué te dijo. En la sesión 9 empieza el proyecto integrador. Ahí vas a usar todo lo del año sobre un problema real de tu entorno.
\end{infoband}

\end{document}
